\documentclass[11pt,letterpaper]{article}

\usepackage[margin=1in]{geometry}
\usepackage{amsmath,amssymb,amsthm,mathtools}
\usepackage{enumitem}
\usepackage{booktabs}
\usepackage{hyperref}
\usepackage{microtype}

\numberwithin{equation}{section}

\theoremstyle{plain}
\newtheorem{theorem}{Theorem}[section]
\newtheorem{lemma}[theorem]{Lemma}
\newtheorem{proposition}[theorem]{Proposition}
\newtheorem{corollary}[theorem]{Corollary}
\theoremstyle{definition}
\newtheorem{definition}[theorem]{Definition}
\newtheorem{example}[theorem]{Example}
\theoremstyle{remark}
\newtheorem{remark}[theorem]{Remark}

\newcommand{\NN}{\mathbb{N}}
\newcommand{\ZZ}{\mathbb{Z}}
\DeclareMathOperator{\lcmop}{lcm}
\DeclareMathOperator{\Atom}{Atom}
\DeclareMathOperator{\widthop}{width}

\title{Upper fibers and witness covers for the Fibonacci order-of-apparition map}
\author{
Haobo Ma\\
AELF PTE LTD.\\
\#14-02, Marina Bay Financial Centre Tower 1, 8 Marina Blvd\\
Singapore 018981, Singapore\\
\texttt{auric@aelf.io}
\and
Wenlin Zhang\thanks{Correspondence: \texttt{e1327962@u.nus.edu}}\\
National University of Singapore\\
Singapore\\
\texttt{e1327962@u.nus.edu}
}
\date{}

\begin{document}
\maketitle

\begin{abstract}
For each positive integer $n$, let $B_n$ denote the set of divisors of the
Fibonacci number $F_n$ whose rank of apparition equals~$n$. We prove that
$B_n$ is an upper set in the divisor lattice of $F_n$ obtained by removing
the covering ideals associated with the maximal proper divisors of~$n$. Our
main result classifies the minimal generators of~$B_n$: every minimal
generator factors uniquely as a pairwise coprime product of atomic
prime-power birth moduli forming a witness cover whose support spans the prime
coordinates of~$n$. This yields the bound $\omega(m)\le\omega(n)$ for every
minimal generator $m$ and a primitive-versus-ladder dichotomy for atomic
factors. We give complete classifications for $\omega(n)\le 3$ and show that
both $B_n$ and its minimal-generator set admit unique factorizations into
connected coordinate blocks, producing partition-lattice formulas for $A(n) =
\#B_n$ and $\#\mathcal{M}_n$. As consequences, nonexceptional birth layers
contain canonical principal intervals, Boolean subposets of rank $\tau(n) -
O(1)$, and satisfy Bell-type lower bounds on odd layers. Numerical data for
$n\le 30$ illustrate these structures.
\end{abstract}

\medskip
\noindent\textbf{Keywords.} Fibonacci numbers, rank of apparition, divisor
lattices, upper sets, witness covers, atomic decompositions.

\smallskip
\noindent\textbf{Mathematics Subject Classification (2020).}
11B39, 11A07, 11A25.

%% ====================================================================
\section{Introduction}\label{sec:intro}
%% ====================================================================

For an integer $q\ge 2$, the \emph{rank of apparition} (or \emph{order of
appearance}) $\alpha(q)$ is the least positive index $n$ such that $q\mid
F_n$, where $(F_n)_{n\ge 0}$ denotes the Fibonacci sequence with $F_0=0$,
$F_1=1$. For each $n\ge 2$, write
\[
D_n:=\operatorname{Div}(F_n),
\qquad
B_n:=\{q\in D_n:\ \alpha(q)=n\},
\qquad
A(n):=\#B_n.
\]
Thus $B_n$ consists of those moduli that first divide a Fibonacci number at
index~$n$. The counting function $A(n)$ is subsumed by the
$\alpha$-contraction formalism of Stroi\'nski \cite{Stroinski2016AlphaContraction};
see also the classical M\"obius inversions discussed in
\cite{Halton1966DivisibilityFibonacci,FultonMorris1969ArithmeticalFunctions}.
In this paper we go beyond enumeration and study $B_n$ as an
ordered subset of the divisor lattice $(D_n,\mid)$, asking how it sits inside
$D_n$, how it is generated, and what combinatorial structure it carries.

Three classical facts underlie the analysis. First, $q\mid F_m$ if and only if
$\alpha(q)\mid m$
\cite{Wall1960FibonacciModm,Vinson1963RankApparition}. Second, the Fibonacci
sequence satisfies strong divisibility: $\gcd(F_a,F_b)=F_{\gcd(a,b)}$.
Third, ranks of apparition respect least common multiples:
$\alpha(\lcmop(u,v))=\lcmop(\alpha(u),\alpha(v))$. These properties, combined
with the primitive divisor theorem for Fibonacci numbers
\cite{Carmichael1913PrimitiveDivisors,BiluHanrotVoutier2001PrimitiveDivisorsLucas},
are the sole inputs for the main results.

This point of view complements several nearby strands in the literature.
Luca and Pomerance \cite{LucaPomerance2014LocalBehavior} study the local
behavior of $\alpha$ under additive shifts; Trojovsk\'a
\cite{Trojovska2020PeriodicPoints} and FitzGibbons--Javaheri--Miller--Verga
\cite{FitzGibbonsJavaheriMillerVerga2025Dynamics} study periodic points and
iteration dynamics of $\alpha$; Cameron--Counts--Lund--Miller--Piechnik--Wong
\cite{CameronCountsLundMillerPiechnikWong2022Fibotomic} study fibotomic
polynomials and related algebraic factorizations; and Sanna
\cite{Sanna2016LucasSequences} and Marques \cite{Marques2012OrderAppearance}
study the growth and fixed-point behavior of the order-of-appearance map. The
present paper instead fixes a single fiber $B_n$ and studies its divisibility
geometry inside $D_n$. We are not aware of earlier work that formulates these
fibers as upper sets in divisor lattices or that gives the two-support and
three-support classifications proved here.

Our first main result identifies $B_n$ as an upper fiber obtained by sieve.

\begin{theorem}\label{thm:intro-upper-fiber}
For every $n\ge 2$, $B_n=D_n\setminus\bigcup_{p\mid n}D_{n/p}$, where the
union runs over primes $p\mid n$. In particular, $B_n$ is an upper set in
$D_n$ whose minimal elements $\mathcal M_n:=\min(B_n)$ determine $B_n$
through the interval decomposition $B_n=\bigcup_{m\in\mathcal M_n}[m,F_n]$.
\end{theorem}

Call a prime power $\theta=p^e$ \emph{atomic for the index $d$} if
$\alpha(\theta)=d$ and $\alpha(p^{e-1})<d$. An \emph{$n$-witness cover} is a
pairwise coprime set of such atomic prime powers whose apparition indices
collectively span the prime factorization of $n$ in a minimal, irredundant
fashion (see Definition~\ref{def:witness-cover} for the precise formulation).

\begin{theorem}\label{thm:intro-witness}
For every $n\ge 2$, the map $\Theta\mapsto\prod_{\theta\in\Theta}\theta$ is
a bijection from the set of $n$-witness covers to $\mathcal M_n$.
Consequently, every $m\in\mathcal M_n$ satisfies $\omega(m)\le\omega(n)$.
\end{theorem}

Making this classification effective requires the prime-power apparition
formulas of Fulton--Morris \cite{FultonMorris1969ArithmeticalFunctions} and
Lengyel \cite{Lengyel1995OrderFibonacciLucas}. These show that every atomic
factor is either \emph{primitive} (arising from a new prime) or lies on a
\emph{prime-power ladder} (changing exactly one prime coordinate of the
apparition index). This dichotomy
(Theorem~\ref{thm:primitive-ladder-dichotomy}) enables complete if-and-only-if
classifications of $\mathcal M_n$ when $\omega(n)\le 3$
(Corollary~\ref{cor:two-support-complete-classification} and
Theorem~\ref{thm:fibonacci-three-coordinate-types}).

Our third main result provides a factorization theory.

\begin{theorem}\label{thm:intro-connected}
For every $q\in B_n$, there exist a unique partition $\Pi\in
\operatorname{Part}(\mathcal P(n))$ and unique connected elements
$q_C\in B_{n_C}$ ($C\in\Pi$) such that $q=\prod_{C\in\Pi}q_C$,
where $n_C=\prod_{\ell\in C}\ell^{\nu_\ell(n)}$ and ``connected'' refers to
the support hypergraph of Definition~\ref{def:support-hypergraph}. The
analogous statement holds for $\mathcal M_n$. Consequently,
\[
A(n)=\sum_{\Pi\in\operatorname{Part}(\mathcal P(n))}
\prod_{C\in\Pi}A^{\mathrm{conn}}(n_C),
\]
with a corresponding M\"obius-inversion formula for
$A^{\mathrm{conn}}(n)$.
\end{theorem}

As secondary consequences, nonexceptional birth layers contain canonical
principal intervals
(Corollary~\ref{thm:canonical-principal-interval}), Boolean subposets of rank
$\tau(n)-O(1)$ (Corollary~\ref{thm:boolean-skeleton}), and---for odd
$n$---Bell-type lower bounds on $\#\mathcal M_n$
(Corollary~\ref{cor:odd-primitive-cover-realization}).

The paper is organized as follows. Section~\ref{sec:prelim} collects the
background on Fibonacci divisibility and the prime-power apparition formulas.
Section~\ref{sec:upper-fibers} develops the upper-fiber and witness-cover
theory and proves Theorems~\ref{thm:intro-upper-fiber}
and~\ref{thm:intro-witness}. Section~\ref{sec:low-support} treats the
primitive--ladder dichotomy and the low-support classifications.
Section~\ref{sec:connected} establishes the connected-block factorization
(Theorem~\ref{thm:intro-connected}) and the partition-lattice formulas.
Section~\ref{sec:numerical} presents numerical data, and
Section~\ref{sec:questions} records further questions.

%% ====================================================================
\section{Preliminaries}\label{sec:prelim}
%% ====================================================================

\subsection{Rank of apparition and Fibonacci divisibility}\label{subsec:fib-div}

\begin{proposition}\label{prop:rank-basics}
Let $q\ge 2$ and $n\ge 1$.
\begin{enumerate}[label=\textup{(\alph*)}]
\item One has $q\mid F_n$ if and only if $\alpha(q)\mid n$.
\item For any finite family $q_1,\dots,q_s\ge 2$,
$\alpha(\lcmop(q_1,\dots,q_s))=\lcmop(\alpha(q_1),\dots,\alpha(q_s))$.
\item For all $a,b\ge 1$, $\gcd(F_a,F_b)=F_{\gcd(a,b)}$ and
$F_a\mid F_b$ if and only if $a\mid b$.
\end{enumerate}
\end{proposition}

\begin{proof}
Parts (a) and (c) are due to Wall \cite{Wall1960FibonacciModm}; see also
Vinson \cite{Vinson1963RankApparition}. Part~(b) follows from (a): writing
$Q=\lcmop(q_1,\dots,q_s)$, one has $Q\mid F_n$ if and only if
$q_i\mid F_n$ for all $i$, which holds if and only if
$\lcmop(\alpha(q_1),\dots,\alpha(q_s))\mid n$.
\end{proof}

\subsection{Primitive divisors}\label{subsec:primitive}

\begin{theorem}\label{thm:primitive-divisors}
For every $n\ge 3$ with $n\notin\{6,12\}$, there exists a prime $p_n$ with
$\alpha(p_n)=n$.
\end{theorem}

\begin{proof}
For $n=3,4,5$ one takes $p_3=2$, $p_4=3$, $p_5=5$. For $n\ge 7$, $n\neq 12$,
the primitive divisor theorem
\cite{Carmichael1913PrimitiveDivisors,Yabuta2001PrimitiveDivisorFibonacci,%
BiluHanrotVoutier2001PrimitiveDivisorsLucas}
gives a prime dividing $F_n$ but no earlier Fibonacci number. Since
$\alpha(p_n)\mid n$ and $p_n$ divides no $F_d$ with $d<n$, necessarily
$\alpha(p_n)=n$.
\end{proof}

\subsection{Counting corollaries}\label{subsec:counting}

The standard M\"obius-inversion formulas for first-appearance counts follow
from Stroi\'nski's $\alpha$-contraction identity
\cite{Stroinski2016AlphaContraction}: for any $f:\NN\to\mathbb{C}$ with
$f(1)=0$,
\begin{equation}\label{eq:stroinski}
\sum_{d\mid F_n}f(d)=\sum_{d\mid n}f_\alpha(d),
\qquad
f_\alpha(d):=\sum_{\alpha(m)=d}f(m).
\end{equation}
Taking $f$ to be the indicator of integers at least $2$, of primes, or of
prime powers, and applying M\"obius inversion, one obtains
\begin{equation}\label{eq:mobius-birth}
A(n)=\sum_{d\mid n}\mu(n/d)\,\tau(F_d),
\qquad
A_{\mathrm{pr}}(n)=\sum_{d\mid n}\mu(n/d)\,\omega(F_d).
\end{equation}
We regard these as known corollaries; the results of
Sections~\ref{sec:upper-fibers}--\ref{sec:connected} are independent of them.

\subsection{Prime-power apparition formulas}\label{subsec:prime-power}

\begin{proposition}\label{prop:prime-power-ladders}
Let $p$ be prime and $r\ge 1$.
\begin{enumerate}[label=\textup{(\alph*)}]
\item $\alpha(5^r)=5^r$.
\item $\alpha(2)=3$, $\alpha(4)=\alpha(8)=6$, and
$\alpha(2^r)=3\cdot 2^{r-2}$ for $r\ge 3$.
\item If $p\neq 2,5$ and $e_p:=\nu_p(F_{\alpha(p)})$, then
$\alpha(p^r)=\alpha(p)\,p^{\max(0,r-e_p)}$.
\end{enumerate}
\end{proposition}

\begin{proof}
Part~(a) follows from $\nu_5(F_n)=\nu_5(n)$. For odd $p\neq 5$, part~(c) is
proved in \cite{FultonMorris1969ArithmeticalFunctions} using Lengyel's
valuation formula \cite{Lengyel1995OrderFibonacciLucas}:
$\nu_p(F_n)=\nu_p(n)+e_p$ when $\alpha(p)\mid n$. The formula for powers of
$2$ appears in Jarden \cite{Jarden1966RecurringSequences}.
\end{proof}

\begin{definition}\label{def:atomic-families}
For each $d\ge 2$, define the \emph{atomic birth family}
\[
\Atom(d):=\{p^e:\ \alpha(p^e)=d,\ \alpha(p^{e-1})<d\}.
\]
\end{definition}

\begin{corollary}\label{cor:atomic-classification}
Every power $5^r$ belongs to $\Atom(5^r)$. For $p=2$, the atomic powers are
$2\in\Atom(3)$, $4\in\Atom(6)$, and $2^r\in\Atom(3\cdot 2^{r-2})$ for
$r\ge 4$. For $p\neq 2,5$ with $z_p:=\alpha(p)$ and $e_p:=\nu_p(F_{z_p})$,
one has $p\in\Atom(z_p)$ and $p^r\in\Atom(z_p\,p^{r-e_p})$ for $r\ge e_p+1$;
no power $p^r$ with $2\le r\le e_p$ is atomic.
\end{corollary}

\begin{proof}
This is a direct translation of
Proposition~\ref{prop:prime-power-ladders}.
\end{proof}

%% ====================================================================
\section{Birth layers and witness covers}\label{sec:upper-fibers}
%% ====================================================================

\begin{definition}\label{def:birth-layer}
For each $n\ge 2$, let $D_n:=\operatorname{Div}(F_n)$, and define the
\emph{birth layer} $B_n:=\{q\in D_n:\ \alpha(q)=n\}$. We write
$\widthop(B_n)$ for the maximum cardinality of an antichain in the poset
$(B_n,\mid)$.
\end{definition}

\begin{theorem}\label{thm:upper-fiber}
For every $n\ge 2$,
\[
B_n=D_n\setminus\bigcup_{p\mid n}D_{n/p}.
\]
\end{theorem}

\begin{proof}
If $q\in B_n$, then $\alpha(q)=n$, so $q\nmid F_{n/p}$ for every prime
$p\mid n$; otherwise $n=\alpha(q)\mid n/p<n$, a contradiction. Conversely, if
$q\in D_n\setminus\bigcup D_{n/p}$, then $d:=\alpha(q)\mid n$. Were $d<n$,
some prime $p$ would divide $n/d$, giving $d\mid n/p$ and $q\mid F_{n/p}$, a
contradiction. Hence $\alpha(q)=n$.
\end{proof}

\begin{theorem}\label{thm:minimal-generators}
The set $\mathcal M_n:=\min(B_n)$ satisfies:
\begin{enumerate}[label=\textup{(\alph*)}]
\item $B_n$ is an upper set in $D_n$ with unique maximal element $F_n$;
\item $B_n=\bigcup_{m\in\mathcal M_n}[m,F_n]$, where
$[m,F_n]:=\{d\in D_n:\ m\mid d\}$;
\item for $q\in D_n$, one has $q\in\mathcal M_n$ if and only if
$\alpha(q)=n$ and $\alpha(q/p)<n$ for every prime $p\mid q$.
\end{enumerate}
If $q=\prod_{i=1}^r p_i^{e_i}$ and $\lambda_i:=\alpha(p_i^{e_i})$,
$\lambda_i^-:=\alpha(p_i^{e_i-1})$, then $q\in\mathcal M_n$ if and only if
\begin{equation}\label{eq:lcm-criterion}
\lcmop(\lambda_1,\dots,\lambda_r)=n
\end{equation}
and, for every $i$,
$\lcmop(\lambda_1,\dots,\lambda_{i-1},\lambda_i^-,\lambda_{i+1},\dots,
\lambda_r)<n$.
\end{theorem}

\begin{proof}
If $q\in B_n$ and $q\mid d\mid F_n$, then $\alpha(d)\mid n$ and
$\alpha(q)\mid\alpha(d)$; since $\alpha(q)=n$, one gets $\alpha(d)=n$, so
$d\in B_n$. Thus $B_n$ is an upper set. Since $\alpha(F_n)=n$
(Proposition~\ref{prop:rank-basics}(c)), $F_n$ is the unique maximum.
Every element lies above a minimal element, giving~(b).

For~(c), if $q\in\mathcal M_n$ and $p\mid q$, then $q/p\notin B_n$ by
minimality, so $\alpha(q/p)<n$. Conversely, if $\alpha(q)=n$ and every
$\alpha(q/p)<n$, then any $d\mid q$ with $d\in B_n$ and $d<q$ would yield
$q/p\in B_n$ for some prime $p\mid q/d$, a contradiction. The lcm
reformulation follows from
Proposition~\ref{prop:rank-basics}(b).
\end{proof}

\noindent\emph{Proof of Theorem~\ref{thm:intro-upper-fiber}.}
This follows from Theorems~\ref{thm:upper-fiber} and
\ref{thm:minimal-generators}.\qed

\subsection*{Witness covers}

\begin{definition}\label{def:witness-cover}
Let $n=\prod_{\ell\in\mathcal P(n)}\ell^{a_\ell}$. For
$\theta=p^e\in\Atom(d)$ with $d\mid n$, write $d(\theta):=\alpha(\theta)=d$,
$d^-(\theta):=\alpha(p^{e-1})$, and define
\[
T_n(\theta):=\{\ell\in\mathcal P(n):\ \nu_\ell(d(\theta))=a_\ell\},
\qquad
E_n(\theta):=\{\ell\in\mathcal P(n):\ \nu_\ell(d(\theta))=a_\ell>
\nu_\ell(d^-(\theta))\}.
\]
A finite set $\Theta\subseteq\bigcup_{d\mid n}\Atom(d)$ is an
\emph{$n$-witness cover} if its elements are pairwise coprime,
$\bigcup_{\theta\in\Theta}T_n(\theta)=\mathcal P(n)$, and, for every
$\theta\in\Theta$,
$E_n(\theta)\nsubseteq\bigcup_{\eta\neq\theta}T_n(\eta)$.
\end{definition}

\begin{theorem}\label{thm:witness-classification}
For every $n\ge 2$,
$\mathcal M_n=\bigl\{\prod_{\theta\in\Theta}\theta:\
\Theta\text{ is an $n$-witness cover}\bigr\}$.
\end{theorem}

\begin{proof}
Let $m\in\mathcal M_n$, $m=\prod_{i=1}^r\theta_i$ ($\theta_i=p_i^{e_i}$).
Each $\theta_i$ is atomic: if $\alpha(p_i^{e_i-1})=\alpha(p_i^{e_i})$, then
replacing $\theta_i$ by $p_i^{e_i-1}$ preserves the lcm, contradicting
minimality via Theorem~\ref{thm:minimal-generators}(c). Since
$n=\lcmop(d(\theta_i))$, the sets $T_n(\theta_i)$ cover $\mathcal P(n)$. The
condition $\alpha(m/p_i)<n$ forces
$E_n(\theta_i)\nsubseteq\bigcup_{j\neq i}T_n(\theta_j)$; otherwise the lost
coordinates would be supplied by other factors, keeping the lcm equal to $n$.
Thus $\{\theta_1,\dots,\theta_r\}$ is an $n$-witness cover.

Conversely, let $\Theta$ be an $n$-witness cover and
$m=\prod_{\theta\in\Theta}\theta$. By
Proposition~\ref{prop:rank-basics}(b), $\alpha(m)=\lcmop(d(\theta):
\theta\in\Theta)=n$. For each $\theta=p^e\in\Theta$, the witness condition
provides $\ell\in E_n(\theta)\setminus\bigcup_{\eta\neq\theta}T_n(\eta)$; at
this coordinate, $\alpha(m/p)<n$.
Theorem~\ref{thm:minimal-generators}(c) then gives $m\in\mathcal M_n$.
\end{proof}

\begin{corollary}\label{cor:support-bound}
Every $m\in\mathcal M_n$ satisfies $\omega(m)\le\omega(n)$.
\end{corollary}

\begin{proof}
The witness condition provides an injection
$\theta\mapsto\ell(\theta)$ from $\Theta$ into $\mathcal P(n)$: choose
$\ell(\theta)\in E_n(\theta)\setminus\bigcup_{\eta\neq\theta}T_n(\eta)$.
Hence $\omega(m)=\#\Theta\le\omega(n)$.
\end{proof}

\noindent\emph{Proof of Theorem~\ref{thm:intro-witness}.}
This follows from Theorem~\ref{thm:witness-classification} and
Corollary~\ref{cor:support-bound}.\qed

\begin{example}\label{ex:n20}
For $n=20$, $F_{20}=6765=3\cdot 5\cdot 11\cdot 41$. The relevant atomic
factors are $3\in\Atom(4)$, $5\in\Atom(5)$, $11\in\Atom(10)$,
$41\in\Atom(20)$. The witness covers are $\{41\}$, $\{3,5\}$, $\{3,11\}$,
yielding $\mathcal M_{20}=\{15,33,41\}$ and $A(20)=11$.
\end{example}

\begin{example}\label{ex:n30}
For $n=30$, $F_{30}=832040=2^3\cdot 5\cdot 11\cdot 31\cdot 61$. The witness
covers are $\{31\}$, $\{4,5\}$, $\{2,11\}$, $\{4,61\}$, $\{11,61\}$, giving
$\mathcal M_{30}=\{20,22,31,244,671\}$ and $A(30)=52$.
\end{example}

\subsection*{Substructures of nonexceptional birth layers}

\begin{corollary}\label{thm:canonical-principal-interval}
Let $n\ge 3$ with $n\notin\{6,12\}$, and let $\pi_n$ be a primitive prime
with $\alpha(\pi_n)=n$ (Theorem~\ref{thm:primitive-divisors}). Then
$\pi_n\in\mathcal M_n$, and
\[
A(n)\ge\tau(F_n/\pi_n)\ge\tfrac{1}{2}\,\tau(F_n).
\]
\end{corollary}

\begin{proof}
The singleton $\{\pi_n\}$ is an $n$-witness cover with
$E_n(\pi_n)=T_n(\pi_n)=\mathcal P(n)$, so $\pi_n\in\mathcal M_n$. The
interval $[\pi_n,F_n]$ has $\tau(F_n/\pi_n)$ elements, and
$\tau(F_n/\pi_n)=\tau(F_n)\cdot\nu_{\pi_n}(F_n)/(\nu_{\pi_n}(F_n)+1)
\ge\tau(F_n)/2$.
\end{proof}

\begin{corollary}\label{thm:boolean-skeleton}
Let $n\ge 3$ with $n\notin\{6,12\}$, let
$\mathcal D^\sharp(n):=\{d\mid n:\ d\notin\{1,2,6,12\}\}$, and for each
$d\in\mathcal D^\sharp(n)$ choose a prime $\pi_d$ with $\alpha(\pi_d)=d$.
Set $r(n):=\#(\mathcal D^\sharp(n)\setminus\{n\})$. Then the map
$S\mapsto\pi_n\prod_{d\in S}\pi_d$ embeds the Boolean lattice $B_{r(n)}$
into $B_n$. In particular,
\[
A(n)\ge 2^{r(n)},
\qquad
\widthop(B_n)\ge\binom{r(n)}{\lfloor r(n)/2\rfloor}.
\]
\end{corollary}

\begin{proof}
For each such $S$, the product has apparition index
$\lcmop(n,d:d\in S)=n$, hence lies in $B_n$. The primes are pairwise
distinct, so divisibility is equivalent to set inclusion. The width bound
follows from Sperner's theorem \cite{Sperner1928Subsets}.
\end{proof}

\begin{theorem}\label{thm:maximal-order}
One has
\[
\limsup_{n\to\infty}
\frac{\log\log A(n)\cdot\log\log n}{\log n}\ge\log 2.
\]
The same bound holds with $\widthop(B_n)$ in place of $A(n)$.
\end{theorem}

\begin{proof}
By Corollary~\ref{thm:boolean-skeleton}, $\log_2 A(n)\ge r(n)=\tau(n)+O(1)$.
Wigert's theorem \cite{Wigert1907Divisors} gives
$\limsup(\log\tau(n)\cdot\log\log n/\log n)=\log 2$,
and the result follows.
\end{proof}

%% ====================================================================
\section{Low-support classifications}\label{sec:low-support}
%% ====================================================================

\begin{theorem}\label{thm:primitive-ladder-dichotomy}
Let $\theta=p^e$ be an atomic birth modulus in some $n$-witness cover. Then
exactly one of the following holds:
\begin{enumerate}[label=\textup{(\alph*)}]
\item $d^-(\theta)=1$ and $E_n(\theta)=T_n(\theta)$
\textup{(primitive atom)};
\item $d^-(\theta)>1$ and $\#E_n(\theta)=1$
\textup{(ladder atom)}.
\end{enumerate}
\end{theorem}

\begin{proof}
If $d^-(\theta)=1$, then by definition
$E_n(\theta)=T_n(\theta)$.
If $d^-(\theta)>1$, Corollary~\ref{cor:atomic-classification} shows that in
every case ($p=2$, $p=5$, or $p\neq 2,5$), lowering the exponent changes
exactly one prime coordinate of the apparition index, so $E_n(\theta)$ is a
singleton.
\end{proof}

\begin{theorem}\label{thm:extremal-coordinate-separation}
Let $k:=\omega(n)$ and $m\in\mathcal M_n$. Then $\omega(m)=k$ if and only
if, after indexing $\mathcal P(n)=\{\ell_1,\dots,\ell_k\}$, one can write
$m=\theta_1\cdots\theta_k$ with pairwise coprime atomic birth moduli
satisfying $T_n(\theta_i)=\{\ell_i\}$ for each~$i$.
\end{theorem}

\begin{proof}
If $\omega(m)=k$, the injection $\theta\mapsto\ell(\theta)$ of
Corollary~\ref{cor:support-bound} is bijective, and each $T_n(\theta)$ is
forced to be a singleton. Conversely, if each $T_n(\theta_i)=\{\ell_i\}$,
then $E_n(\theta_i)=\{\ell_i\}$, the conditions of
Definition~\ref{def:witness-cover} are immediate, and
Theorem~\ref{thm:witness-classification} gives $m\in\mathcal M_n$ with
$\omega(m)=k$.
\end{proof}

\begin{corollary}\label{cor:two-support-complete-classification}
If $\mathcal P(n)=\{\ell_1,\ell_2\}$, define
$\Atom_{\ell_i}(n):=\{\theta\in\bigcup_{d\mid n}\Atom(d):
T_n(\theta)=\{\ell_i\}\}$. Then
\[
\mathcal M_n
=\Atom(n)\sqcup
\{xy:\ x\in\Atom_{\ell_1}(n),\ y\in\Atom_{\ell_2}(n),\ \gcd(x,y)=1\}.
\]
\end{corollary}

\begin{proof}
Every non-atomic $m\in\mathcal M_n$ has
$\omega(m)=2=\omega(n)$ by Corollary~\ref{cor:support-bound}.
Theorem~\ref{thm:extremal-coordinate-separation} then gives the stated
decomposition.
\end{proof}

We now classify the three-support case. The abstract enumeration is a finite
exercise; the arithmetic content lies in identifying which abstract types
survive the primitive--ladder constraint.

\begin{theorem}\label{thm:abstract-three-types}
Let $\Theta$ be an abstract witness family on $[3]$, meaning a finite family
of pairs $(E,T)$ with $\varnothing\neq E\subseteq T\subseteq[3]$,
$\bigcup T=[3]$, and
$E\nsubseteq\bigcup_{(E',T')\neq(E,T)}T'$ for each $(E,T)$. Up to
relabeling $[3]$, exactly nine types occur:
\begin{align*}
&\Gamma_1=\{(123,123)\},\quad
\Gamma_2=\{(12,123)\},\quad
\Gamma_3=\{(1,123)\},\\
&\Gamma_4=\{(12,12),(23,23)\},\quad
\Gamma_5=\{(1,12),(23,23)\},\quad
\Gamma_6=\{(1,12),(3,23)\},\\
&\Gamma_7=\{(1,1),(23,23)\},\quad
\Gamma_8=\{(1,1),(2,23)\},\quad
\Gamma_9=\{(1,1),(2,2),(3,3)\}.
\end{align*}
\end{theorem}

\begin{proof}
The witness condition injects $\Theta$ into $[3]$, so $\#\Theta\le 3$.
If $\#\Theta=1$, then $T=[3]$ and $E$ has size $3$, $2$, or $1$, giving
$\Gamma_1$, $\Gamma_2$, $\Gamma_3$.
If $\#\Theta=3$, the injection is bijective and forces every $T$ to be a
singleton, giving $\Gamma_9$.
If $\#\Theta=2$, write $\Theta=\{(E_1,T_1),(E_2,T_2)\}$. Neither $T_i$ is
$[3]$ (else the witness condition for the other factor fails), so each has
size $1$ or $2$. Since $T_1\cup T_2=[3]$, there are two support patterns up
to relabeling: $(T_1,T_2)=(\{1\},\{2,3\})$ yielding $\Gamma_7$ and
$\Gamma_8$, and $(T_1,T_2)=(\{1,2\},\{2,3\})$ yielding $\Gamma_4$,
$\Gamma_5$, $\Gamma_6$.
\end{proof}

\begin{theorem}\label{thm:fibonacci-three-coordinate-types}
If $\omega(n)=3$, then every witness cover of an element of $\mathcal M_n$
belongs to one of the eight types
$\Gamma_1,\Gamma_3,\Gamma_4,\Gamma_5,\Gamma_6,\Gamma_7,\Gamma_8,\Gamma_9$.
The type $\Gamma_2$ is excluded: it requires a factor with
$1<\#E<\#T=3$, which is forbidden by the primitive--ladder dichotomy
\textup{(Theorem~\ref{thm:primitive-ladder-dichotomy})}.
\end{theorem}

\begin{proof}
By Theorem~\ref{thm:primitive-ladder-dichotomy}, every Fibonacci atomic
factor satisfies either $E=T$ or $\#E=1$. The unique abstract type with
$1<\#E<\#T$ is $\Gamma_2$.
\end{proof}

\begin{corollary}\label{cor:three-coordinate-realization}
If $\omega(n)=3$, every element of $\mathcal M_n$ arises by making pairwise
coprime selections from one of the eight families corresponding to
$\Gamma_1,\Gamma_3,\Gamma_4,\dots,\Gamma_9$ (see
Theorem~\ref{thm:fibonacci-three-coordinate-types}), where each family is
determined by: the primitive-atom families
$\mathcal P_S(n):=\{\theta:E_n(\theta)=T_n(\theta)=S\}$ and the ladder-atom
families
$\mathcal L_{i,S}(n):=\{\theta:E_n(\theta)=\{i\},T_n(\theta)=S\}$ for
nonempty $S\subseteq[3]$ and $i\in S$.
\end{corollary}

\begin{remark}\label{rem:wss}
The Wall--Sun--Sun condition enters through
Proposition~\ref{prop:prime-power-ladders}(c): for $p\neq 2,5$, one has
$\alpha(p^2)=\alpha(p)$ if and only if $\nu_p(F_{\alpha(p)})\ge 2$, which is
precisely the definition of a Wall--Sun--Sun prime
\cite{SunSun1992FibonacciFLT}. No such primes are currently known, and their
nonexistence would imply that every odd prime contributes an immediate ladder
step above the first atomic level.
\end{remark}

%% ====================================================================
\section{Connected block factorizations}\label{sec:connected}
%% ====================================================================

Throughout this section, write
$n_C:=\prod_{\ell\in C}\ell^{\nu_\ell(n)}$ for nonempty
$C\subseteq\mathcal P(n)$.

\begin{definition}\label{def:support-hypergraph}
If $q=\prod_i p_i^{e_i}\in B_n$, its \emph{support hypergraph}
$\mathcal H_n(q)$ has vertex set $\mathcal P(n)$ and edge set
$\{\mathcal P(\alpha(p_i^{e_i})):p_i^{e_i}\parallel q\}$. We call $q$
\emph{connected} if $\mathcal H_n(q)$ is connected, and write
$A^{\mathrm{conn}}(n)$ for the number of connected elements of $B_n$.
Analogous notation applies to minimal generators.
\end{definition}

\begin{theorem}\label{thm:coprime-embedding}
Let $n_1,\dots,n_r\ge 2$ be pairwise coprime and $n=\prod n_i$. Then the
multiplication map $\prod B_{n_i}\to B_n$ and its restriction
$\prod\mathcal M_{n_i}\to\mathcal M_n$ are order embeddings. Consequently,
$A(n)\ge\prod A(n_i)$, $\#\mathcal M_n\ge\prod\#\mathcal M_{n_i}$, and
$\widthop(B_n)\ge\prod\widthop(B_{n_i})$.
\end{theorem}

\begin{proof}
Strong divisibility gives $\gcd(F_{n_i},F_{n_j})=1$ for $i\neq j$, so the
factors are coprime and Proposition~\ref{prop:rank-basics}(b) yields
$\alpha(\prod q_i)=\lcmop(n_i)=n$. The order-embedding property follows from
the cross-index coprimeness. For the minimal-generator embedding, if some
$m_i\notin\mathcal M_{n_i}$, then $\alpha(m_i/p)=n_i$ for some prime
$p$, which by the disjointness of the prime-coordinate supports would give
$\alpha(m/p)=n$, contradicting $m\in\mathcal M_n$.
\end{proof}

\begin{theorem}\label{thm:birth-connected-factorization}
For every $q\in B_n$, there exist a unique partition
$\Pi(q)\in\operatorname{Part}(\mathcal P(n))$ and unique connected elements
$q_C\in B_{n_C}$ ($C\in\Pi(q)$) such that $q=\prod_{C\in\Pi(q)}q_C$.
Conversely, if $\Pi\in\operatorname{Part}(\mathcal P(n))$ and
$q_C\in B_{n_C}$ is connected for every $C\in\Pi$, then
$\prod q_C\in B_n$ with connected decomposition~$\Pi$. In particular,
\begin{equation}\label{eq:partition-formula}
A(n)=\sum_{\Pi\in\operatorname{Part}(\mathcal P(n))}
\prod_{C\in\Pi}A^{\mathrm{conn}}(n_C),
\end{equation}
and the M\"obius inversion on the partition lattice gives
\begin{equation}\label{eq:connected-inversion}
A^{\mathrm{conn}}(n)=\sum_{\Pi\in\operatorname{Part}(\mathcal P(n))}
(-1)^{|\Pi|-1}(|\Pi|-1)!\prod_{C\in\Pi}A(n_C).
\end{equation}
\end{theorem}

\begin{proof}
Write $q=\prod_i\varpi_i$ as its prime-power factorization with
$\lambda_i:=\alpha(\varpi_i)$. Let $\Pi(q)$ be the partition of
$\mathcal P(n)$ into the connected components of $\mathcal H_n(q)$, and for
each block $C$ set $q_C:=\prod_{\mathcal P(\lambda_i)\subseteq C}\varpi_i$.
The factors are pairwise coprime, and
Proposition~\ref{prop:rank-basics}(b) gives $\alpha(q_C)=n_C$ for each~$C$:
every prime coordinate of $n_C$ is realized at full exponent by some factor
of~$q_C$, because $\alpha(q)=n$.

For the converse, Theorem~\ref{thm:coprime-embedding} gives
$\prod q_C\in B_n$, and the edge structure ensures the connected components
are exactly the blocks of~$\Pi$. The counting formula follows by summation
over all partitions, and \eqref{eq:connected-inversion} is the standard
M\"obius inversion on the partition lattice.
\end{proof}

\begin{theorem}\label{thm:mingenerator-connected-factorization}
The analogous statement holds for $\mathcal M_n$: every $m\in\mathcal M_n$
factors uniquely as $m=\prod_{C\in\Pi_M(m)}m_C$ with each $m_C$
a connected minimal generator in $\mathcal M_{n_C}$, and
\[
\#\mathcal M_n=\sum_{\Pi\in\operatorname{Part}(\mathcal P(n))}
\prod_{C\in\Pi}M^{\mathrm{conn}}(n_C).
\]
\end{theorem}

\begin{proof}
Given $m\in\mathcal M_n$, its birth-layer factorization
$m=\prod m_C$ from Theorem~\ref{thm:birth-connected-factorization}
satisfies $m_C\in B_{n_C}$. If some $m_C\notin\mathcal M_{n_C}$, then
there exists a prime $p\mid m_C$ with $\alpha(m_C/p)=n_C$, which by the
disjointness of coordinate supports gives $\alpha(m/p)=n$, contradicting
$m\in\mathcal M_n$. The counting formula follows as before.
\end{proof}

\noindent\emph{Proof of Theorem~\ref{thm:intro-connected}.}
This follows from Theorems~\ref{thm:birth-connected-factorization}
and~\ref{thm:mingenerator-connected-factorization}.\qed

\begin{corollary}\label{cor:four-support-reduction}
If $\omega(n)=4$, every nonconnected element of $\mathcal M_n$ factors
uniquely along a nontrivial partition of $\mathcal P(n)$ with block types
$3{+}1$, $2{+}2$, $2{+}1{+}1$, or $1{+}1{+}1{+}1$, using connected
minimal generators on each block. In particular, the only genuinely new
four-coordinate content lies in the connected cores.
\end{corollary}

\begin{corollary}\label{cor:odd-primitive-cover-realization}
If $n$ is odd, then for every irreducible cover $\mathcal U$ of
$\mathcal P(n)$ by nonempty subsets, the product
$m_{\mathcal U}:=\prod_{C\in\mathcal U}\pi_C$ (where $\alpha(\pi_C)=n_C$) is
a minimal generator. In particular,
\[
\#\mathcal M_n\ge\mathrm{Bell}(\omega(n)),
\qquad
\widthop(B_n)\ge\mathrm{Bell}(\omega(n)).
\]
\end{corollary}

\begin{proof}
Odd $n$ ensures $n_C\notin\{6,12\}$ for every $C$, so
Theorem~\ref{thm:primitive-divisors} provides the primes $\pi_C$. Each
$\pi_C$ is a primitive atom with $E_n(\pi_C)=T_n(\pi_C)=C$, so the
irreducible-cover condition matches the witness-cover condition. The map
$\mathcal U\mapsto m_{\mathcal U}$ is injective, and the number of set
partitions of $\mathcal P(n)$ is $\mathrm{Bell}(\omega(n))$.
\end{proof}

%% ====================================================================
\section{Numerical evidence}\label{sec:numerical}
%% ====================================================================

Table~\ref{tab:birth-layer-data} displays the birth layer cardinality $A(n)$
and the number of minimal generators $\#\mathcal M_n$ for $2\le n\le 30$,
computed by exhaustive enumeration. Several features visible in the data are
consistent with the structural results of the preceding sections.
First, $\#\mathcal M_n=1$ whenever $\omega(n)=1$ and
$n\notin\{19,25,27\}$; the exceptions arise from composite Fibonacci numbers
with multiple atomic families at the same prime-power index.
Second, the jump at $n=12$ ($A(12)=10$, $\#\mathcal M_{12}=3$) and at $n=24$
($A(24)=55$, $\#\mathcal M_{24}=4$) reflects the interaction of the
exceptional prime $2$ with the non-primitive layers at indices $6$ and $12$.
Third, the first three-coordinate layer at $n=30$ already exhibits
$\#\mathcal M_{30}=5$ witness covers spanning all eight admissible types from
Theorem~\ref{thm:fibonacci-three-coordinate-types}.

\begin{table}[ht]
\centering
\small
\caption{Birth layer data for $2\le n\le 30$.}
\label{tab:birth-layer-data}
\begin{tabular}{rrcrrl}
\toprule
$n$ & $\omega(n)$ & $F_n$ & $A(n)$ & $\#\mathcal{M}_n$ &
$\mathcal{M}_n$ \\
\midrule
 3 & 1 & $2$ & 1 & 1 & $\{2\}$\\
 4 & 1 & $3$ & 1 & 1 & $\{3\}$\\
 5 & 1 & $5$ & 1 & 1 & $\{5\}$\\
 6 & 2 & $2^3$ & 2 & 1 & $\{4\}$\\
 7 & 1 & $13$ & 1 & 1 & $\{13\}$\\
 8 & 1 & $3\cdot 7$ & 2 & 1 & $\{7\}$\\
 9 & 1 & $2\cdot 17$ & 2 & 1 & $\{17\}$\\
 10 & 2 & $5\cdot 11$ & 2 & 1 & $\{11\}$\\
 11 & 1 & $89$ & 1 & 1 & $\{89\}$\\
 12 & 2 & $2^4\cdot 3^2$ & 10 & 3 & $\{6,9,16\}$\\
 13 & 1 & $233$ & 1 & 1 & $\{233\}$\\
 14 & 2 & $13\cdot 29$ & 2 & 1 & $\{29\}$\\
 15 & 2 & $2\cdot 5\cdot 61$ & 5 & 2 & $\{10,61\}$\\
 16 & 1 & $3\cdot 7\cdot 47$ & 4 & 1 & $\{47\}$\\
 17 & 1 & $1597$ & 1 & 1 & $\{1597\}$\\
 18 & 2 & $2^3\cdot 17\cdot 19$ & 10 & 2 & $\{19,68\}$\\
 19 & 1 & $37\cdot 113$ & 3 & 2 & $\{37,113\}$\\
 20 & 2 & $3\cdot 5\cdot 11\cdot 41$ & 11 & 3 & $\{15,33,41\}$\\
 21 & 2 & $2\cdot 13\cdot 421$ & 5 & 2 & $\{26,421\}$\\
 22 & 2 & $89\cdot 199$ & 2 & 1 & $\{199\}$\\
 23 & 1 & $28657$ & 1 & 1 & $\{28657\}$\\
 24 & 2 & $2^5\cdot 3^2\cdot 7\cdot 23$ & 55 & 4 & $\{14,23,32,63\}$\\
 25 & 1 & $5^2\cdot 3001$ & 4 & 2 & $\{25,3001\}$\\
 26 & 2 & $233\cdot 521$ & 2 & 1 & $\{521\}$\\
 27 & 1 & $2\cdot 17\cdot 53\cdot 109$ & 12 & 2 & $\{53,109\}$\\
 28 & 2 & $3\cdot 13\cdot 29\cdot 281$ & 11 & 3 & $\{39,87,281\}$\\
 29 & 1 & $514229$ & 1 & 1 & $\{514229\}$\\
 30 & 3 & $2^3\cdot 5\cdot 11\cdot 31\cdot 61$ & 52 & 5 &
$\{20,22,31,244,671\}$\\
\bottomrule
\end{tabular}
\end{table}

%% ====================================================================
\section{Further questions}\label{sec:questions}
%% ====================================================================

\begin{enumerate}[label=(\arabic*)]
\item Extend the three-support classification to sharp arithmetic
nonemptiness criteria for each family in
Corollary~\ref{cor:three-coordinate-realization}, then resolve the connected
cores when $\omega(n)=4$. By
Corollary~\ref{cor:four-support-reduction}, only the connected part is
genuinely new.
\item Obtain arithmetic or asymptotic control of the connected counts
$A^{\mathrm{conn}}(n)$ and $M^{\mathrm{conn}}(n)$. The partition-lattice
formulas reduce global invariants to these kernels but do not yet estimate
their frequency or size.
\item Extend the witness-cover formalism to Lucas sequences and, more
generally, to strong divisibility sequences equipped with a
rank-of-apparition map and a primitive-divisor theorem.
\end{enumerate}

%% ====================================================================
\section*{Declarations}
%% ====================================================================

\noindent\textbf{Funding.} No specific funding was received for this work.

\smallskip
\noindent\textbf{Author contributions.} Both authors contributed to the
conception of the project, the development of the arguments, and the writing
of the manuscript. Both authors read and approved the final version.

\smallskip
\noindent\textbf{Conflict of interest.} The authors declare that they have no
conflict of interest.

\smallskip
\noindent\textbf{Data availability.} The computational data in
Table~\ref{tab:birth-layer-data} were produced by a Python script available
from the authors upon request.

\bibliographystyle{plain}
\bibliography{references_local}

\end{document}
